% SPDX-License-Identifier: LPPL-1.3c
% Copyright (C) 2026 Christophe Jorssen
% Author and Current Maintainer: Christophe Jorssen <christophe.jorssen@gmail.com>
% LPPL maintenance status: maintained
% This file is part of luafemm. See LICENSE and LICENSES.md for its terms.
\section{The dielectric catalogue}
\label{sec:dielectrics}

The electric interface includes the 95 entries in FEMM~5's electrostatic
|statlib.dat|, independently of the magnetic material catalogue. These are
nominal properties, with the material conditions retained in the original
names below. No network access or native FEMM installation is needed when
compiling a document. Data provenance, the pinned upstream revision and byte
checksums are supplied in |data/femm5/README.md|. The unchanged data and their
transcription retain the MIT notice in |data/femm5/LICENSE.txt|; the original
package implementation remains LPPL. FEMM~4 dielectric data are not bundled.

\begin{key}{/femm/electrostatic/material=\meta{custom or catalogue identifier} (initially custom)}
Picture-local choice of background dielectric for analytical or FE electric
models. This is a PGF |.is choice| handler: use an identifier from the table,
not the original material name. Selection assigns its relative permittivity.
It does not add a region, source or conducting electrode. Background volume
charge remains an explicit |volume charge density| input for a compatible
model. All catalogue entries currently have zero tabulated volume charge.

Ordinary key order applies. A following |relative permittivity| overrides
the selected scalar value; selecting the material again resets it. |custom|
clears the catalogue identifier and retains the current numeric settings.
Selecting any material while restoring a snapshot is a physical override and
is rejected. The snapshot already owns its material and constitutive values.
\end{key}
\begin{key}{/femm/electric region/material=\meta{custom or catalogue identifier} (initially custom)}
Path-local FE dielectric choice, available in both planar and axisymmetric
models. Selection assigns |relative permittivity| and |volume charge density|
from the record; later explicit keys override either value. |custom| retains
those numeric settings while clearing the identifier. The existing region
fill rule, overlap order, local mesh size and curve tolerance are unchanged.
A style can collect this key with the usual native drawing options.
\end{key}

\begin{codeexample}[width=0pt]
\tikzset{electric glass/.style={fill=blue!10,
    femm/electric region={material=glass-soda-lime}}}
\begin{tikzpicture}[femm/electrostatic={model=finite elements,
    xmin=0,xmax=2,ymin=0,ymax=1,outer boundary=insulated,
    material=air-dry-co2-free,mesh size=.2}]
  \draw[electric glass] (0,0) rectangle (1,1);
  \path[femm/electric region={material=ptfe,relative permittivity=2.2}]
    (1,0) rectangle (2,1);
  \draw[thick,femm/electric electrode={potential=10}] (0,0)--(0,1);
  \draw[thick,femm/electric electrode={potential=0}] (2,0)--(2,1);
  \pic {femm electric field};\pic {femm equipotentials};
  \node[below,font=\scriptsize] at (.5,0) {Glass};
  \node[below,font=\scriptsize] at (1.5,0) {PTFE};
\end{tikzpicture}
\end{codeexample}

There are 94 isotropic records and one anisotropic record, silicon carbide
with its crystal axis along $y$. All 95 are usable by the FE provider. Material
selection installs both principal values at angle zero; subsequent scalar
|relative permittivity| assignments deliberately replace them with an isotropic
value. Analytical providers reject the anisotropic catalogue entry.
The tensor keys and their precedence are documented in
Section~\ref{sec:electric-extended}. In axisymmetry the directions are $r,z$.
No dispersion, dielectric loss, temperature dependence or nonlinear dielectric
law is inferred from a name.

The following table is generated from the same pinned source as the runtime
choices. It lists relative permittivities (dimensionless); every volume charge
density in this edition is zero. Use catalogue values as starting points,
and use measured or supplier-certified properties when an application needs
them. Trade names identify the data and imply no endorsement.

\input{\luafemmdocdir/dielectrics-table.tex}
